\section{Discussion and Limitations}
The design intentionally separates the measurement problem from the modelling problem. A low-channel representation can appear robust because preprocessing inadvertently uses target-subject statistics, because participants leak across windows, or because different corpus labels are mapped to superficially similar categories. Each failure invalidates a broad generalization claim. Recent evidence of severe identity confounding in qEEG benchmarks makes these checks indispensable \cite{pandilova2026}.

The benchmark cannot establish physiological equivalence between different hardware, and affective ratings are operational reference annotations rather than direct neural ground truth. Potentially small eligible corpus overlap may prevent a cross-dataset factorial test; that is a scientifically interpretable feasibility outcome, not a reason to interpolate incompatible sensors. A negative or near-chance external result is publishable only insofar as the protocol is sufficiently rigorous, novel and transparent.

\section{Conclusion}
The proposed study is intended to measure the limits of sparse EEG transfer under explicitly audited montage, person and domain changes. Completion depends on legitimate dataset access, frozen analysis decisions and reproducible baselines. No clinical or diagnostic inference is proposed.
